% --- INSTITUTIONAL TEMPLATE FED v7.4.9 ---
% Estándar de Reporteo APA 7 (Babel-Free)
\documentclass[12pt,spanish]{article}

\usepackage{graphicx}
\usepackage{geometry}
\usepackage{indentfirst} 
\usepackage{caption}
\usepackage{floatrow} 

% Forzar estilos de caption antes de cualquier carga
\floatsetup[figure]{capposition=top}
\floatsetup[table]{capposition=top}

\usepackage{booktabs}
\usepackage{float}
\usepackage{longtable}
\usepackage{threeparttable}
\usepackage{array}
\usepackage{calc}
\usepackage{amsmath, amssymb}
\usepackage{fancyhdr}
\usepackage{mathptmx} 
\usepackage[T1]{fontenc}
\usepackage{setspace} 
\usepackage{ragged2e} 
\usepackage{titlesec}
\usepackage{url}
\def\UrlBreaks{\do\/\do-\do\.\do\_\do\?\do\=\do\&\do\%}
\usepackage[hidelinks,breaklinks]{hyperref}

% --- Pandoc Compatibility Block (Iron Clad) ---
\providecommand{\tightlist}{%
  \setlength{\itemsep}{0pt}\setlength{\parskip}{0pt}}

\usepackage{xcolor}
\usepackage{fancyvrb}
\usepackage{framed}

% --- Code Highlighting (Pandoc Standard) ---
\AtBeginDocument{
  \definecolor{shadecolor}{RGB}{248,248,248}
}

% Compatibilidad Pandoc 3.x
\ifdefined\pandocbounded\else
  \newcommand{\pandocbounded}[1]{#1}
\fi

% --- Pandoc Citeproc (CSL) compatibility ---
\newlength{\cslhangindent}
\setlength{\cslhangindent}{1.27cm} 
\newlength{\csllabelwidth}
\setlength{\csllabelwidth}{3em}
\newcommand{\citeproctext}{} 
\newcommand{\citeprocitem}[2]{\item} 
\newenvironment{CSLReferences}[2] 
 {\begin{list}{}{
  \setlength{\leftmargin}{\cslhangindent}
  \setlength{\parsep}{0pt}
  \setlength{\itemsep}{6pt}
  \ifodd #1
   \setlength{\leftmargin}{\leftmargin + \csllabelwidth}
   \setlength{\itemindent}{-\csllabelwidth}
  \fi
  }
  \renewcommand{\bibitem}[2][]{\item} % Elimina los corchetes [ ]
 }
 {\end{list}}
\newcommand{\CSLBlock}[1]{#1\hfill\break}
\newcommand{\CSLLeftMargin}[1]{\parbox[t]{\csllabelwidth}{#1}}
\newcommand{\CSLRightInline}[1]{\parbox[t]{\linewidth - \csllabelwidth}{#1}\break}
\newcommand{\CSLIndent}[1]{\hspace{\cslhangindent}#1}

\makeatletter
\@ifundefined{paragraph}{\newcounter{paragraph}}{}
\@ifundefined{subparagraph}{\newcounter{subparagraph}}{}
\makeatother

% Solución al error 'No counter none defined' en longtable
\makeatletter
\newcounter{none}
\@ifpackageloaded{caption}{
  \DeclareCaptionType{none}
  \captionsetup[none]{name=Tabla}
}{}
\makeatother

\hypersetup{
  colorlinks=true,
  linkcolor=blue,
  filecolor=blue,
  citecolor=blue,
  urlcolor=blue
}

% --- Macros Institucionales ---
\newcommand{\fuente}[1]{
  \vspace{4pt}
  \begin{flushleft}
    \small
    \textit{Nota.} #1
  \end{flushleft}
}

% --- Localización Manual Determinista (Sin Babel) ---
\AtBeginDocument{
  \renewcommand{\tablename}{Tabla}
  \renewcommand{\figurename}{Figura}
  \renewcommand{\contentsname}{Índice}
  \renewcommand{\listtablename}{Índice de Tablas}
  \renewcommand{\listfigurename}{Índice de Figuras}
  \renewcommand{\abstractname}{Resumen}
  \renewcommand{\refname}{Referencias}
  
  \RaggedRight
  \setlength{\parindent}{1.27cm}
  \setlength{\parskip}{0pt}
}

\DeclareCaptionLabelSeparator{newline}{\\ \vspace{2pt}}
\captionsetup{
  position=top,
  justification=RaggedRight,
  singlelinecheck=false,
  labelfont=bf,
  textfont=it,
  labelsep=newline,
  font=small,
  skip=10pt
}
\captionsetup[figure]{name=Figura}
\captionsetup[table]{name=Tabla}

% Márgenes APA 7
\geometry{
  a4paper,
  margin=25.4mm,
  headheight=15mm
}

% --- Encabezados ---
\pagestyle{fancy}
\fancyhf{} 
\fancyhead[L]{\includegraphics[height=1.2cm]{\logoUNL}}
\fancyhead[R]{\includegraphics[height=1.2cm]{\logoECONOMIA}}
\renewcommand{\headrulewidth}{0pt}
\fancyfoot[R]{\thepage}

\fancypagestyle{plain}{
  \fancyhf{}
  \fancyhead[L]{\includegraphics[height=1.2cm]{\logoUNL}}
  \fancyhead[R]{\includegraphics[height=1.2cm]{\logoECONOMIA}}
  \fancyfoot[R]{\thepage}
  \renewcommand{\headrulewidth}{0pt}
}

% --- Títulos ---
\titleformat{\section}{\normalfont\normalsize\bfseries\centering}{\thesection}{1em}{}
\titleformat{\subsection}{\normalfont\normalsize\bfseries}{\thesubsection}{1em}{}
\titleformat{\subsubsection}{\normalfont\normalsize\bfseries\itshape}{\thesubsubsection}{1em}{}
\titleformat{\paragraph}{\normalfont\normalsize\bfseries}{\theparagraph}{1em}{}
\titlespacing*{\section}{0pt}{12pt}{6pt}
\titlespacing*{\subsection}{0pt}{12pt}{6pt}
\titlespacing*{\subsubsection}{0pt}{12pt}{6pt}
\titlespacing*{\paragraph}{0pt}{12pt}{6pt}

\newenvironment{referenciasAPA}{
  \setlength{\parindent}{-1.27cm}
  \setlength{\leftskip}{1.27cm}
  \setlength{\parskip}{6pt}
}{\par}

\usepackage{fontspec}
\newcommand{\logoUNL}{/home/erick-fcs/Capital_Workstation/capital-workstation-libs/src/ecs_quantitative/reporting/assets/logos/logo_unl.png}
\newcommand{\logoECONOMIA}{/home/erick-fcs/Capital_Workstation/capital-workstation-libs/src/ecs_quantitative/reporting/assets/logos/logo_economia.jpg}
\renewcommand{\RaggedRight}{\justifying}\makeatletter\@ifpackageloaded{caption}{}{\usepackage{caption}}\AtBeginDocument{%
\ifdefined\contentsname
  \renewcommand*\contentsname{Table of contents}
\else
  \newcommand\contentsname{Table of contents}
\fi
\ifdefined\listfigurename
  \renewcommand*\listfigurename{List of Figures}
\else
  \newcommand\listfigurename{List of Figures}
\fi
\ifdefined\listtablename
  \renewcommand*\listtablename{List of Tables}
\else
  \newcommand\listtablename{List of Tables}
\fi
\ifdefined\figurename
  \renewcommand*\figurename{Figure}
\else
  \newcommand\figurename{Figure}
\fi
\ifdefined\tablename
  \renewcommand*\tablename{Table}
\else
  \newcommand\tablename{Table}
\fi
}\@ifpackageloaded{float}{}{\usepackage{float}}
\floatstyle{ruled}
\@ifundefined{c@chapter}{\newfloat{codelisting}{h}{lop}}{\newfloat{codelisting}{h}{lop}[chapter]}
\floatname{codelisting}{Listing}\newcommand*\listoflistings{\listof{codelisting}{List of Listings}}\makeatother\makeatletter\makeatother\makeatletter\@ifpackageloaded{caption}{}{\usepackage{caption}}
\@ifpackageloaded{subcaption}{}{\usepackage{subcaption}}\makeatother

\begin{document}

\begin{center}
    {\bfseries \Large UNIVERSIDAD NACIONAL DE LOJA}\\[0.3cm]
    {\bfseries \large FACULTAD JURÍDICA, SOCIAL Y
ADMINISTRATIVA}\\[0.2cm]
    {\bfseries CARRERA DE ECONOMÍA}\\[1.5cm]
    
    {\bfseries \large POLÍTICA ECONÓMICA}\\[1cm]
    
    {\bfseries \Large Organizador Gráfico de la Política
Monetaria}\\[1.5cm]
\end{center}

\begin{center}
    \setstretch{1.1}
                        Erick FCS \\
                \vspace{0.5cm}
    Econ. Maria Gabriela Moreno Hurtado \\ 
    11 de julio de 2026
\end{center}

\vspace{1cm}

\doublespacing
\setlength{\parindent}{1.27cm}

Este documento presenta el organizador gráfico detallado del
\textbf{Capítulo 11: Política Monetaria} estructurado como una
\textbf{Ruta Causal e Integración Conceptual} que conecta los
instrumentos, canales de transmisión, limitaciones y objetivos de la
política monetaria.

\begin{center}\rule{0.5\linewidth}{0.5pt}\end{center}

\section{1. Diagrama de la Ruta Causal
(Mermaid)}\label{diagrama-de-la-ruta-causal-mermaid}

El siguiente diagrama de flujo ilustra el canal de transmisión de las
decisiones monetarias desde los instrumentos iniciales hasta los
objetivos finales del Banco Central.

\includegraphics[width=16.7in,height=18.76in]{organizador_files/figure-latex/mermaid-figure-1.png}

\begin{center}\rule{0.5\linewidth}{0.5pt}\end{center}

\section{2. Componentes Conceptuales del
Organizador}\label{componentes-conceptuales-del-organizador}

\subsection{2.1. Concepto y Rol de la Política
Monetaria}\label{concepto-y-rol-de-la-poluxedtica-monetaria}

\begin{itemize}
\tightlist
\item
  \textbf{Definición:} Regulación de la cantidad de dinero, el crédito y
  las tasas de interés domésticas administrada por la autoridad
  monetaria (Banco Central).
\item
  \textbf{Influencia:} Afecta directamente las reservas del sistema
  bancario interbancario e indirectamente las expectativas y la
  producción agregada.
\end{itemize}

\subsection{2.2. Objetivos Principales}\label{objetivos-principales}

\begin{itemize}
\tightlist
\item
  \textbf{Estabilidad de Precios:} Preservar el valor nominal de la
  moneda mediante la contención de shocks de precios.
\item
  \textbf{Estabilidad de Crecimiento y Empleo:} Mitigar las
  fluctuaciones recesivas del PIB estimulando la demanda interna.
\item
  \textbf{Estabilidad Cambiaria:} Regular la liquidez en divisas
  internacionales para sostener la solvencia de pagos con el exterior.
\end{itemize}

\subsection{2.3. Autonomía de la Autoridad
Monetaria}\label{autonomuxeda-de-la-autoridad-monetaria}

\begin{itemize}
\tightlist
\item
  \textbf{Propósito:} Proteger la gestión técnica del dinero de
  presiones fiscales electorales oportunistas.
\item
  \textbf{Evidencia:} Reduce la inconsistencia temporal e impide el
  financiamiento directo de déficits presupuestarios con emisión de
  dinero inorgánico.
\end{itemize}

\subsection{2.4. Estrategias Operativas}\label{estrategias-operativas}

\begin{itemize}
\tightlist
\item
  \emph{Agregados Monetarios:} Control rígido de la oferta monetaria
  bajo el supuesto de velocidad de circulación estable del dinero.
\item
  \emph{Metas de Inflación:} Ajuste preventivo del tipo de interés
  interbancario de referencia para guiar las expectativas de precios de
  mediano plazo.
\end{itemize}

\subsection{2.5. Instrumentos Causales y Efectos
Esperados}\label{instrumentos-causales-y-efectos-esperados}

{\def\LTcaptype{none} % do not increment counter
\begin{longtable}[]{@{}
  >{\raggedright\arraybackslash}p{(\linewidth - 6\tabcolsep) * \real{0.2500}}
  >{\raggedright\arraybackslash}p{(\linewidth - 6\tabcolsep) * \real{0.2500}}
  >{\raggedright\arraybackslash}p{(\linewidth - 6\tabcolsep) * \real{0.2500}}
  >{\raggedright\arraybackslash}p{(\linewidth - 6\tabcolsep) * \real{0.2500}}@{}}
\toprule\noalign{}
\begin{minipage}[b]{\linewidth}\raggedright
Instrumento
\end{minipage} & \begin{minipage}[b]{\linewidth}\raggedright
Medida
\end{minipage} & \begin{minipage}[b]{\linewidth}\raggedright
Efecto Intermedio
\end{minipage} & \begin{minipage}[b]{\linewidth}\raggedright
Efecto en Economía Real
\end{minipage} \\
\midrule\noalign{}
\endhead
\bottomrule\noalign{}
\endlastfoot
\textbf{OMA (Venta)} & Venta de títulos a la banca & \(\downarrow\)
Reservas líquidas de bancos & \(\uparrow\) Tasas de interés de mercado
\(\rightarrow\) \(\downarrow\) Consumo e Inversión \(\rightarrow\)
\(\downarrow\) Demanda Agregada. \\
\textbf{Encaje Bancario} & \(\uparrow\) Coeficiente legal &
\(\downarrow\) Capacidad crediticia del sistema & \(\uparrow\) Costo de
financiamiento \(\rightarrow\) \(\downarrow\) Gasto privado agregado
\(\rightarrow\) \(\downarrow\) Inflación. \\
\textbf{Redescuento} & \(\uparrow\) Tasa de redescuento & \(\downarrow\)
Incentivo a solicitar liquidez & Actitud restrictiva bancaria
\(\rightarrow\) Menor canalización de préstamos. \\
\end{longtable}
}

\subsection{2.6. Limitaciones
Estructurales}\label{limitaciones-estructurales}

\begin{itemize}
\tightlist
\item
  \emph{Asimetría de impacto:} Elevada efectividad en fases de
  sobrecalentamiento (política contractiva) y nula efectividad en
  recesiones profundas (política expansiva) si hay desconfianza y
  aversión al crédito.
\item
  \emph{Retardos de implementación:} Los efectos finales sobre las
  variables reales demoran entre 6 y 18 meses tras la toma de decisión.
\end{itemize}

\begin{center}\rule{0.5\linewidth}{0.5pt}\end{center}

\section{3. Recuadro Especial: El Caso de Ecuador (Dolarización y
Gobierno de Daniel
Noboa)}\label{recuadro-especial-el-caso-de-ecuador-dolarizaciuxf3n-y-gobierno-de-daniel-noboa}

La dolarización oficial adoptada en el año 2000 eliminó la soberanía
monetaria del país. Durante el gobierno de \textbf{Daniel Noboa
(2024-2026)}, el funcionamiento monetario se ha caracterizado por las
siguientes dinámicas empíricas de gran impacto:

\subsection{3.1. Protección de las Reservas y Acuerdo con el FMI
(2024)}\label{protecciuxf3n-de-las-reservas-y-acuerdo-con-el-fmi-2024}

La debilidad de la caja fiscal presionó el balance del Banco Central del
Ecuador (BCE) por el uso de depósitos del sector público. En mayo de
2024, la administración de Daniel Noboa concretó un acuerdo de Servicio
Ampliado del Fondo (SAF) con el FMI por 4000 millones de dólares. El
ingreso de estos recursos externos blindó la liquidez de las reservas
internacionales del BCE, evitando una descapitalización que pusiera en
riesgo el esquema de dolarización.

\subsection{3.2. Regulación de Techos de Tasas de Interés y Contracción
del
Crédito}\label{regulaciuxf3n-de-techos-de-tasas-de-interuxe9s-y-contracciuxf3n-del-cruxe9dito}

Ante la contracción del crédito de la banca privada provocada por el
incremento del riesgo país en 2024, el BCE reformó la metodología para
fijar los techos máximos de las tasas de interés de los microcréditos y
segmentos productivos. El objetivo fue reabrir el canal de crédito,
flexibilizando las tasas para que la banca pudiera cubrir los costos de
intermediación en entornos de alta incertidumbre.

\subsection{3.3. Control del Coeficiente de Liquidez Doméstica
(LDR)}\label{control-del-coeficiente-de-liquidez-domuxe9stica-ldr}

Para contrarrestar la fuga de capitales y la contracción de la liquidez
física en el país, el BCE mantuvo regulados los porcentajes mínimos de
liquidez que la banca privada debe conservar dentro del territorio
nacional (liquidez doméstica). Esta medida mitigó la salida de divisas
hacia bancos del exterior, obligando al sistema financiero a recircular
el crédito de manera interna bajo el ciclo económico de Noboa.

\begin{center}\rule{0.5\linewidth}{0.5pt}\end{center}

\section{4. Explicación Breve Escrita a Mano (Síntesis de una
Página)}\label{explicaciuxf3n-breve-escrita-a-mano-suxedntesis-de-una-puxe1gina}

Respuestas conceptuales ultra-sintetizadas para transcripción a mano en
una página:

\begin{itemize}
\tightlist
\item
  \textbf{1. Idea central del Capítulo 11:} La política monetaria es un
  instrumento macroeconómico indirecto que gestiona la liquidez y las
  tasas de interés mediante el sistema financiero para estabilizar
  precios y empleo. Su efectividad depende de la credibilidad y
  autonomía del Banco Central para coordinar las expectativas de los
  agentes y evitar retardos internos o externos en su transmisión.
\item
  \textbf{2. Relación entre instrumentos, canales y objetivos:} El Banco
  Central activa instrumentos (OMA, encaje legal o redescuento) que
  alteran las reservas bancarias primarias y el tipo de interés
  interbancaria. Este shock se propaga por canales de transmisión
  (crédito, tasas, activos y expectativas), modificando el costo de
  financiamiento y ajustando la demanda agregada para alcanzar los
  objetivos finales de baja inflación y crecimiento.
\item
  \textbf{3. Efectos distintos en el corto y largo plazo:} En el corto
  plazo, debido a la rigidez temporal de precios y salarios, la masa
  monetaria influye sobre variables reales como la demanda, la
  producción y el empleo (no neutralidad). En el largo plazo, los
  precios se ajustan completamente, disipando el impacto real y
  generando solo un aumento proporcional del nivel de precios
  (neutralidad del dinero), haciendo que la inflación sea un fenómeno
  puramente monetario.
\item
  \textbf{4. Limitaciones en economía dolarizada (Ecuador):} Al perder
  la soberanía cambiaria y la emisión primaria de dinero, el país carece
  de prestamista de última instancia incondicional y tipo de cambio
  nominal para absorber shocks externos. La política monetaria nacional
  queda limitada al control micro-prudencial de la liquidez doméstica
  (encaje y coeficientes de liquidez) y a la regulación de techos de
  tasas de interés.
\item
  \textbf{5. Síntesis Final (¿Por qué no es uniforme el análisis?):} La
  política monetaria no puede analizarse igual en todos los países
  porque su efectividad depende del régimen cambiario y del grado de
  soberanía monetaria. Las economías dolarizadas como Ecuador no pueden
  emitir moneda propia ni anclar tasas de referencia activas, limitando
  al Banco Central a custodiar la liquidez doméstica frente a las
  fluctuaciones exógenas de la balanza de pagos.
\end{itemize}

\end{document}
